\chapter{Sensores e dados do veículo}
\label{cap:sensores}

% --- CABEÇA ------------------------------------------------------
Um sistema pericial não decide sobre imagens nem sobre sinais de rádio:
decide sobre factos, como \enquote{há um peão na passadeira} ou \enquote{a distância ao
veículo da frente é insegura}. Num veículo, esses factos nascem de medições através de dados provenientes dos sensores.
Estes sensores não tomam decisões e limitam-se a medir e a detetar dados específicos, e cabe à
base de conhecimento interpretar o que medem e transformá-lo em factos.

Este capítulo acaba por
responder a três perguntas essenciais para o desenvolvimento do sistema:
que sensores fornecem a informação de que o sistema precisa para decidir
(\cref{sec:sensores-considerados}), como é que as suas leituras se
transformam em factos (\cref{sec:sensores-dados}) e em que condições se
pode confiar nesses factos (\cref{sec:sensores-qualidade}). Nesta fase do
projeto, os sensores são simulados, o que permite reproduzir cada cenário
de forma controlada e repetível.

% --- TRONCO 1 ----------------------------------------------------
\section{Sensores considerados}
\label{sec:sensores-considerados}

Os sensores de um veículo autónomo costumam dividir-se em dois grupos: os
exterocetivos, que observam o ambiente em redor, e os propriocetivos, que
medem o estado do próprio veículo, como a velocidade das rodas, as
acelerações ou o ângulo da direção~\cite{yeong2021}. Entre os primeiros,
três sensores assumem um papel fulcral e complementam-se precisamente por
responderem a perguntas diferentes sobre o mesmo objeto: a câmara permite
saber \emph{o que} é (um peão, um sinal, a cor de um semáforo ou uma marca
rodoviária), o radar indica \emph{a que distância} está e \emph{a que
velocidade} se aproxima, e o LiDAR descreve a sua \emph{forma} com maior
precisão. O número de sensores que um veículo utiliza depende, contudo, do
nível de automação, crescendo à medida que o veículo assume mais
responsabilidade pela condução. A \cref{tab:sensores} sintetiza os sensores
considerados para este sistema, indicando, para cada um, a grandeza que
mede, a sua principal vantagem e limitação e a utilização que lhe é dada.

\begin{table}[H]
  \centering
  \caption{Sensores considerados para o sistema}
  \label{tab:sensores}
  \footnotesize
  \renewcommand{\arraystretch}{1.1}
  \begin{tabularx}{\textwidth}{@{}L{2.3cm}V
      >{\raggedright\arraybackslash}XV
      >{\raggedright\arraybackslash}XV
      >{\raggedright\arraybackslash}XV
      >{\raggedright\arraybackslash}X@{}}
    \toprule
    \textbf{Sensor} & \textbf{Grandeza medida} & \textbf{Principal vantagem}
      & \textbf{Principal limitação} & \textbf{Utilização no sistema} \\
    \midrule
    Radar frontal
      & Distância e velocidade relativa
      & Funciona de noite, com chuva e nevoeiro
      & Fraca capacidade de classificação
      & Risco de colisão e distância de segurança \\ \linhasep
    Câmara frontal
      & Tipo de objeto, sinais e marcas rodoviárias
      & Única fonte semântica e de cor
      & Sensível à iluminação e ao nevoeiro
      & Peões, obstáculos, sinalização e saída de via \\ \linhasep
    LiDAR
      & Forma e distância dos objetos em 3D
      & Geometria precisa, sem depender da luz
      & Degrada-se com nevoeiro e chuva; custo
      & Confirmação de obstáculos e desempate \\ \linhasep
    Radares de canto
      & Ocupação das vias adjacentes
      & Cobertura do ângulo morto
      & Dificuldade com objetos parados
      & Mudança de via \\ \linhasep
    Sensor de chuva
      & Intensidade da chuva
      & Baixo custo e uso generalizado
      & Não mede a aderência
      & Velocidade adequada à chuva \\ \linhasep
    Sensor de luminosidade
      & Luz ambiente
      & Deteta a noite e a entrada em túneis
      & Não mede a visibilidade
      & Fiabilidade da câmara \\ \linhasep
    Sensores de velocidade das rodas
      & Velocidade própria e patinagem
      & Robustez
      & Afetados pela patinagem
      & Todas as decisões; aderência \\ \linhasep
    Sensor de ângulo da direção
      & Ângulo do volante
      & Revela a intenção do condutor
      & Insuficiente de forma isolada
      & Saída de via \\
    \bottomrule
  \end{tabularx}
\end{table}

No sistema, estes três sensores deverão trabalhar em conjunto. A câmara e o radar
formam a base, porque se complementam (a câmara reconhece o objeto, mas
estima mal a distância, enquanto o radar mede a distância e a velocidade
relativa com precisão, mesmo de noite ou com chuva, mas distingue mal o tipo
de objeto~\cite{wei2022}) e porque é neles que assentam, na prática, as
funções de segurança hoje obrigatórias nos automóveis novos da União
Europeia, como a travagem automática de emergência~\cite{ue2019_2144}. O
LiDAR completará esta base com a forma precisa de obstáculos e peões e servirá
de terceira fonte quando a câmara e o radar divergem. Perante um obstáculo
à frente, por exemplo, a câmara confirma a sua presença na
trajetória, o radar indica com que rapidez o veículo se aproxima dele e o
LiDAR delimitará a sua posição e dimensão. Os sensores do próprio veículo,
de velocidade das rodas e de ângulo da direção, dão o contexto de que todas
as regras precisam. Ficam de fora os sensores ultrassónicos, cujo alcance de
poucos metros serve sobretudo as manobras de estacionamento, e o conjunto
GNSS e unidade de medição inercial, porque a localização global não é
necessária nas situações tratadas pelo sistema. Nem todos os sensores
aqui descritos estão ainda implementados no protótipo, cujo estado é
resumido no \cref{cap:implementacao}.

Estes sensores fornecem medições, mas as regras precisam de factos. A
secção seguinte mostra essa passagem.

% --- TRONCO 2 ----------------------------------------------------
\section{Medições diretas e dados derivados}
\label{sec:sensores-dados}

A informação chega às regras em três níveis: as \emph{medições diretas},
lidas de um sensor, como a distância ao objeto da frente, a velocidade
própria ou o ângulo da direção; os \emph{dados derivados}, calculados a
partir de várias medições, como o tempo até à colisão, o intervalo
temporal ao veículo da frente ou a aderência estimada; e os \emph{factos
semânticos}, como \enquote{o risco de colisão é crítico} ou \enquote{o piso está
escorregadio}, que resultam da comparação de um dado derivado com um limiar. Esse limiar,
ou a regra que o calcula, é conhecimento do perito, e é
precisamente neste ponto que a aquisição de conhecimento se liga aos
sensores.

Para avaliar a distância ao objeto da frente usam-se dois indicadores,
ambos calculados a partir da distância $d$ medida pelo radar. O TTC indica
quantos segundos faltam para o embate se as velocidades
se mantiverem, e obtém-se dividindo $d$ pela velocidade relativa
$v_{\mathrm{rel}}$, como define o Regulamento UNECE n.º~152~\cite{unr152};
o intervalo temporal (THW) indica quanto tempo o veículo demora a chegar ao
ponto onde o objeto está agora, e obtém-se dividindo $d$ pela velocidade
própria $v$:
\begin{equation}
  \mathit{TTC} = \frac{d}{v_{\mathrm{rel}}} \quad (v_{\mathrm{rel}} > 0),
  \qquad\qquad
  \mathit{THW} = \frac{d}{v}.
  \label{eq:ttc-thw}
\end{equation}
Os dois indicadores são complementares: o THW mede o perigo potencial de
seguir demasiado perto, e o TTC mede o perigo real de uma
aproximação~\cite{vogel2003}, embora o protótipo utilize, por agora,
apenas o TTC. Para a saída involuntária de via está previsto um
indicador análogo, o tempo até transpor a marca rodoviária
(TLC)~\cite{mammar2006}.

Para distinguir uma situação normal de uma situação crítica, o perito
indicou que não basta fixar um valor de TTC: a criticidade deve resultar de
uma equação que relacione a velocidade atual, a distância ao objeto e a
capacidade de travagem, sem nunca considerar arranques ou travagens
instantâneas. O sistema deverá, por isso, comparar a distância
$d$ ao objeto com a distância de que precisa para anular a velocidade
relativa, tendo em conta o tempo de reação $t_r$ e a desaceleração máxima
$a$, que depende da aderência, e admitindo que o objeto mantém a sua
velocidade:
\begin{equation}
  d_{\mathrm{trav}} = v_{\mathrm{rel}}\,t_r + \frac{v_{\mathrm{rel}}^{2}}{2a}.
  \label{eq:dtrav}
\end{equation}
Enquanto $d$ excede claramente $d_{\mathrm{trav}}$, a situação é normal;
quando a margem se aproxima de zero, torna-se crítica, embora ainda seja
possível travar ou mudar de via; e, abaixo de $d_{\mathrm{trav}}$, o
veículo ultrapassa o ponto de não-retorno da travagem, restando o desvio,
se a via adjacente estiver livre, ou a travagem de emergência para reduzir
a gravidade do impacto. O TTC continuará a ser útil, mas passará a ser
comparado com o tempo necessário para reagir e travar, em vez de um valor
fixo. Os critérios resultam assim da sessão com o perito, ao passo que os
valores de $t_r$ e de $a$ serão definidos pela equipa e verificados nos
testes, de modo a garantir que o aviso ao condutor surja pelo menos $0,8$~s antes
de uma travagem de emergência, como exige o Regulamento UNECE
n.º~152~\cite{unr152}. Já o intervalo mínimo de segurança partirá das
recomendações nacionais, que apontam para 3~s e para 4 a 5~s com chuva ou
nevoeiro~\cite{prp2025}.

A mesma equação determina o alcance de que os sensores precisam. O Código
da Estrada exige que a velocidade permita parar no espaço livre e visível à
frente~\cite{codigoestrada}, o que, num sistema automatizado, significa parar
dentro do alcance \emph{efetivo} dos sensores. Como esse alcance diminui
com o nevoeiro e a chuva e a distância de travagem aumenta
com o piso molhado, a velocidade tem de baixar quando a visibilidade ou a
aderência descem.

Estes factos, porém, só são úteis se forem fiáveis, o que leva à qualidade
dos dados.

% --- PERNAS ------------------------------------------------------
\section{Qualidade e fusão de dados}
\label{sec:sensores-qualidade}

Para que essa fiabilidade possa ser avaliada, cada facto de origem sensorial
guarda, além do valor, o sensor que o produziu e o instante da leitura, e
deverá guardar também um nível de confiança, informação que serve ainda para
explicar cada decisão.
Seguindo a indicação do perito, uma perceção só é válida enquanto o objeto
se mantém no alcance dos sensores, e os sinais de trânsito não são
memorizados como objetos, limitando-se a alterar o estado do veículo. A
falta de dados de um sensor tem ainda de ser representada explicitamente,
porque, num sistema de regras, a ausência de um facto seria interpretada
como estrada livre, quando na verdade o estado é \emph{desconhecido}.

A confiança numa deteção aumenta quando esta é confirmada, e o número de
leituras exigido deverá depender da criticidade calculada pela \cref{eq:dtrav}: numa
situação normal, o sistema aguardará uma segunda leitura para excluir valores
anómalos, enquanto numa situação crítica agirá logo à primeira, porque esperar
poderia custar o momento de agir. A confirmação pode também vir de sensores
diferentes, já que a câmara e o radar, sendo complementares por medirem
aspetos diferentes da cena, tornam a deteção mais fiável quando observam o
mesmo objeto~\cite{castanedo2013}. Num sistema pericial, esta fusão faz-se
ao nível da decisão, através de regras que combinam as conclusões de cada
sensor.

Mesmo sem avaria, as condições podem levar os sensores a contradizer-se.
Nesse caso, o LiDAR servirá de terceira fonte e, se também ele não permitir
decidir, aplicar-se-á uma regra de desempate, parando o veículo se o nível de
risco for alto ou crítico. A falha de um sensor conduzirá igualmente à
paragem, sempre de forma controlada, já que nenhuma travagem é instantânea.

% --- FECHO -------------------------------------------------------
Em síntese, o sistema assenta na câmara e no radar, a que se juntará o LiDAR,
complementados pelos sensores do próprio veículo, cujas leituras se
transformam em factos através de critérios definidos com o perito e deverão
transportar a informação necessária para avaliar a sua confiança.
